\section{Introducción}

La valuación de derivados y la inferencia estadística sobre el proceso del activo subyacente están estrechamente relacionadas, pero resuelven problemas distintos. Los retornos históricos informan sobre la dinámica observada bajo una medida física, denotada por $\Pmeasure$. En contraste, la teoría de no arbitraje establece que, en un mercado completo como Black--Scholes--Merton, el valor de un reclamo contingente replicable se obtiene como una esperanza descontada bajo una medida neutral al riesgo $\Qmeasure$ \citep{blackscholes1973,merton1973}. Esta distinción adquiere especial relevancia cuando los parámetros son desconocidos: la incertidumbre estimada bajo $\Pmeasure$ sólo debe trasladarse al precio si el parámetro correspondiente también determina la dinámica de valuación bajo $\Qmeasure$.

Este trabajo estudia esa interfaz en una opción asiática de promedio aritmético. Las opciones asiáticas son derivados dependientes de trayectoria cuyo pago se determina a partir del promedio del activo durante un conjunto de fechas de monitoreo. Son naturales cuando la exposición económica depende de precios promedio, como puede ocurrir en materias primas, divisas o compras recurrentes. Incluso bajo Black--Scholes, el promedio aritmético de precios lognormales correlacionados no posee una distribución cerrada simple. Esta dificultad ha dado lugar a métodos analíticos y numéricos específicos. \citet{kemnavorst1990} utilizan la opción sobre promedio geométrico como variable de control en Monte Carlo; \citet{curran1994} condiciona sobre el promedio geométrico; y \citet{kahale2020} desarrolla métodos multilevel para opciones asiáticas discretamente monitoreadas.

En paralelo, la literatura bayesiana de opciones ha mostrado que la incertidumbre paramétrica puede incorporarse a precios y densidades predictivas. \citet{guidolin2003} estudian precios de opciones bajo aprendizaje bayesiano, mientras \citet{romboutsstentoft2014} integran incertidumbre paramétrica dentro de densidades predictivas neutrales al riesgo en modelos heterocedásticos. Un antecedente importante para este artículo es que estos últimos encuentran que el impacto incremental de la incertidumbre paramétrica puede ser pequeño cuando los parámetros se encuentran suficientemente identificados. Por tanto, integrar una posterior completa no implica necesariamente una mejora sustancial del punto de precio.

La literatura mexicana sobre opciones asiáticas se ha concentrado principalmente en enriquecer la dinámica financiera y desarrollar aplicaciones de valuación. \citet{ortiz2016} incorporan tasas de interés estocásticas de tipo Vasicek y CIR; \citet{matias2019} combinan volatilidad estocástica, control óptimo y simulación Monte Carlo; y \citet{ortizjimenez2026} estudian una aplicación a CEMEXCPO con volatilidad realizada y reversión a la media. El presente trabajo adopta una estrategia complementaria. En lugar de incrementar la complejidad dinámica, mantiene Black--Scholes como laboratorio controlado para aislar el efecto de la incertidumbre paramétrica y distinguir con claridad los objetos estimados bajo $\Pmeasure$ de los parámetros relevantes para la valuación bajo $\Qmeasure$.

Bajo esta especificación, los retornos históricos informan sobre $(\mu,\sigma)$, pero el drift físico $\mu$ desaparece de la ecuación de no arbitraje. La incertidumbre de precio inducida por parámetros históricos se concentra, por tanto, en $\sigma$. La pregunta principal es si integrar la distribución posterior completa de volatilidad cambia materialmente el punto de valuación frente a reglas plug-in. Formalmente, se compara
\begin{equation}
 \widehat C_{FB}=\E\!\left[C^{\Qmeasure}(\sigma)\mid\mathcal D\right]
 \qquad\text{con}\qquad
 C^{\Qmeasure}(\widehat\sigma),
 \label{eq:main_question}
\end{equation}
donde $C^{\Qmeasure}(\sigma)$ es el valor de una call asiática aritmética condicionado en una volatilidad dada.

La comparación se realiza mediante muestreo repetido contra un benchmark externo a cada posterior, $C^{\Qmeasure}(\sigma_0)$, donde $\sigma_0$ es conocida por construcción. Esto evita una evaluación circular en la que la media posterior resulte óptima simplemente porque la función de pérdida se calcula contra las mismas muestras posteriores que la definen.

Las contribuciones son cuatro. Primero, se formaliza un flujo coherente entre inferencia histórica bajo $\Pmeasure$ y valuación bajo $\Qmeasure$ para un payoff dependiente de trayectoria. Segundo, se implementa Metropolis--Hastings en $(\mu,\log\sigma)$ y se propaga únicamente la incertidumbre relevante para el precio de no arbitraje. Tercero, se comparan integración posterior completa, media posterior plug-in, MAP y máxima verosimilitud en 81 escenarios que varían cantidad de información histórica, volatilidad verdadera, moneyness y madurez, con 2,000 réplicas independientes por escenario. Cuarto, se documenta un resultado condicional: full Bayes no domina universalmente en error puntual, pero la diferencia respecto del plug-in de media posterior se comporta como predice la curvatura del funcional de precio. Es prácticamente nula alrededor de at-the-money y con historias largas, y aumenta con posteriors más dispersas, contratos fuera de at-the-money y mayores madureces.

El resto del artículo se organiza de la siguiente forma. La sección 2 revisa la literatura y posiciona la contribución. La sección 3 presenta el modelo, la inferencia y el paso a la medida neutral al riesgo. La sección 4 describe el diseño computacional. La sección 5 presenta los resultados. La sección 6 discute limitaciones y extensiones, y la sección 7 concluye.

\section{Estado del arte y posicionamiento}

\subsection{Valuación de opciones asiáticas}

La dificultad fundamental de una opción asiática aritmética surge de que su payoff depende de una suma de variables lognormales correlacionadas. A diferencia del promedio geométrico, cuya distribución puede manejarse analíticamente bajo Black--Scholes, el promedio aritmético no es lognormal. \citet{kemnavorst1990} explotan la cercanía entre ambos promedios para construir una estrategia de reducción de varianza: el valor analítico de la opción geométrica funciona como variable de control para la simulación del payoff aritmético. Este recurso es especialmente útil aquí porque reduce el error Monte Carlo interno y permite separar con mayor claridad la variación debida a incertidumbre paramétrica.

\citet{curran1994} propone una vía condicionada: parte de la distribución conocida del promedio geométrico y calcula la expectativa del payoff condicionando sobre dicho promedio. La literatura posterior ha extendido el repertorio mediante transformadas, ecuaciones diferenciales parciales, quasi-Monte Carlo, importance sampling y esquemas multilevel. \citet{kahale2020}, por ejemplo, muestra que puede reducirse de forma importante el costo de valuación de opciones discretamente monitoreadas mediante una jerarquía de aproximaciones.

Este artículo no propone un nuevo algoritmo de pricing. El Monte Carlo con variable de control se utiliza deliberadamente como una herramienta madura y transparente. La contribución se concentra en el nivel estadístico: cómo propagar incertidumbre sobre $\sigma$ en el funcional de precio sin confundir la dinámica histórica con la dinámica neutral al riesgo.

\subsection{Incertidumbre bayesiana y precios de opciones}

La incertidumbre de parámetros puede afectar la valuación por varios canales. En modelos de aprendizaje, las creencias de los agentes pueden modificar distribuciones de estado y precios. \citet{guidolin2003} muestran que un mecanismo de aprendizaje bayesiano puede producir patrones de volatilidad implícita y mejorar ciertos ejercicios fuera de muestra. Aunque su modelo difiere del utilizado aquí, el resultado establece que tratar parámetros como objetos inciertos puede modificar sistemáticamente las distribuciones relevantes para valuación.

Una conexión más directa aparece en \citet{romboutsstentoft2014}. Los autores realizan inferencia bayesiana sobre modelos de mezclas normales heterocedásticas, derivan dinámicas neutrales al riesgo y construyen densidades predictivas integrando sobre el espacio paramétrico. Conceptualmente, sustituyen la práctica de condicionar en un estimador único por una distribución predictiva que reconoce incertidumbre de parámetros. Su evidencia también impone una cautela importante: el impacto incremental sobre errores de pricing puede ser pequeño.

El presente artículo traslada esta pregunta a un payoff dependiente de trayectoria bajo una especificación más simple. La simplificación es intencional. En Black--Scholes, la separación entre parámetros de $\Pmeasure$ y $\Qmeasure$ es especialmente clara y permite observar el mecanismo sin introducir primas de riesgo no observadas, saltos o factores latentes.

\subsection{Literatura mexicana y brecha aplicada}

La literatura mexicana proporciona antecedentes relevantes sobre simulación y modelos estocásticos para opciones asiáticas. \citet{ortiz2016} estudian opciones asiáticas bajo tasas de interés Vasicek y CIR y realizan calibración mediante máxima verosimilitud. \citet{matias2019} integran volatilidad estocástica, fundamentos microeconómicos, control óptimo y Monte Carlo. \citet{ortizjimenez2026} utilizan volatilidad realizada y reversión a la media en una aplicación a CEMEXCPO. Estos trabajos muestran que la valuación de opciones asiáticas constituye una línea activa dentro de la literatura financiera aplicada mexicana.

La brecha que aborda este artículo es distinta. En el corpus revisado, la atención se dirige principalmente a especificaciones dinámicas y aplicaciones de valuación, mientras el tratamiento explícito de incertidumbre bayesiana sobre parámetros históricos y su propagación coherente desde $\Pmeasure$ hacia un precio bajo $\Qmeasure$ recibe menor atención. El objetivo no es establecer una prioridad absoluta, sino estudiar esta interfaz en un entorno donde el mecanismo puede aislarse y evaluarse mediante simulación controlada.
